% 原书印刷xiii–xvi / PDF14–17，旧版章节号忠实保留，不能套用第二版目录
\section*{第一版前言}\phantomsection\label{front:preface-first}
\addcontentsline{toc}{section}{第一版前言}
当别人得知你正在写一本教科书时，往往会问两个问题中的一个，或者两个都问. 第一个是：“你为什么要写一本书？”第二个是：“你的书与已有的书有什么不同？”第一个问题相当容易回答：因为你对现有教材并不完全满意. 第二个问题则难得多. 答案无法用几句话说清；为了不使听众厌烦——他们也许只是出于礼貌才问——你会试着给出一个简短而机智的回答. 通常并不奏效.

本书的目的是从概率论的基本原理出发，建立理论统计学（theoretical statistics，有别于数理统计学）. 逻辑的发展、证明、思想与主题等，都通过统计论证逐步展开. 因而，我们从概率基础出发，用具有统计意义的技术、定义与概念发展统计推断理论；这些内容又是此前概念的自然延伸与推论. 开始这项工作时，我们并不知道这种做法会有多成功. 当然，最终的评判留给读者.

本书面向统计学专业，或需要专门学习统计学的相关领域的一年级研究生. 先修要求是一年的微积分. 熟悉一些矩阵运算会有帮助，但不是必需的. 本书可用于两个学期或三个学季的统计学入门课程.

前四章介绍概率论基础，并引入后文必需的许多基本内容. 第5章和第6章是最先讨论统计学的两章. 第5章起到概率与统计之间的过渡作用，已有概率基础的学生可由此开始统计理论课程. 第6章颇有特色，详细讨论充分性、似然与不变性三个统计原则，并说明它们在数据建模中的重要性. 并非所有教师都会详细讲授这一章，但我们强烈建议在此花一些时间. 其中，似然原则与不变性原则得到了详细讨论. 它们与充分性原则，以及这些原则背后的思考方式，都是全面理解统计学的基础.

第7至9章构成统计推断的核心：估计，包括点估计与区间估计，以及假设检验. 这些章节的一个主要特点，是区分寻找适当统计技术的方法与评价这些技术的方法. 理论工作者和实践工作者都关心寻找与评价，但我们认为把二者分开十分重要，因为它们关注的问题不同，遵循的规则也不同. 各章题为“其他考虑”的部分也可能引起读者兴趣. 我们在其中说明，统计推断的规则如何像日常实践中那样适当放宽，而仍能产生有意义的推断. 这些部分介绍的许多技术用于统计咨询，有助于分析实际问题并作出推断.

最后三章可以看作专题，但我们认为，任何人的统计教育都应包含对这些材料的一定了解. 第10章全面介绍决策论，并纳入了我们能够包含的最新内容. 第11章讨论方差分析，包括单因素与随机区组情形，从较简单的处理对比理论出发，建立完整分析的理论. 根据我们的经验，实验工作者最关心由对比作出的推断；运用前面发展出的原则，大多数检验与区间都可以由对比推导. 最后，第12章讨论回归理论，先介绍简单线性回归，再讨论“变量含误差”的回归. 后一个主题十分重要，既展示其自身的用途和固有困难，也揭示普通回归推断的局限.

关于如何用本书安排一年课程，我们给出以下更具体的建议. 本书可以用于两类不同的课程. 第一类可称为“偏重数学”，适合有扎实数学基础的统计学专业学生，即至少学过一年半微积分、一些矩阵代数，也许还学过实分析. 对这些学生，我们建议完整讲授第1至9章，大约需要22周，其余时间可从第10至12章中选取专题，按课程需要安排. 学完前九章之后，最后三章各自的内容都是自足的，可以按任意顺序讲授.

另一类课程“偏重实用”. 这种课程也可作为数学基础较强学生的第一门课程，但主要面向学过一年微积分、可能并非统计学专业的学生. 它强调统计理论中更实用的用途，更关心理解基本统计概念并针对不同情形推导合理的统计程序，而较少关注形式上的最优性研究. 这样的课程必然要省略部分材料，但以下各节可以在一年内讲授：
\begin{center}\begin{tabular}{cl}
章 & 节\\\hline
1 & 全部\\
2 & 2.1，2.2，2.3\\
3 & 3.1，3.2\\
4 & 4.1，4.2，4.3，4.5\\
5 & 5.1，5.2，5.3.1，5.4\\
6 & 6.1.1，6.2.1\\
7 & 7.1，7.2.1，7.2.2，7.2.3，7.3.1，7.3.3，7.4\\
8 & 8.1，8.2.1，8.2.3，8.2.4，8.3.1，8.3.2，8.4\\
9 & 9.1，9.2.1，9.2.2，9.2.4，9.3.1，9.4\\
11 & 11.1，11.2\\
12 & 12.1，12.2
\end{tabular}\end{center}
若时间允许，可以讨论第4.4、5.5、6.1.2、6.1.3和6.1.4节的材料，而不必过分强调细节. 第11.3和12.3节的内容也可以考虑.

习题取自许多来源，数量相当丰富. 我们认为，掌握这些材料的唯一途径也许就是练习，因此提供了充足的练习机会. 我们力求习题形式多样，其中许多题目说明了正文之外的新观点，或补充正文内容. 有些习题甚至取自研究论文. 当你能把自己学生时代还是新研究的论文改编成习题时，难免会觉得自己老了！习题虽然不像章节那样细分，但其排列大体遵循章节顺序，因为分组往往会给出过多提示. 此外，随着题号增大，习题的挑战性大体上也会增加.

本书是一本范围较广的入门教材，因此对各主题的讨论并不十分深入. 但我们认为，有责任引导读者在感兴趣的主题上再向前一步. 因此，书中收入了大量参考文献，指出深入理解各个主题的途径. Kotz、Johnson和Read主编的《统计科学百科全书》（\textit{Encyclopedia of Statistical Sciences}）为许多主题提供了很好的入门介绍.

写作本书时，我们借鉴了过去的教学经验与当前的研究工作，也得到了许多人的帮助，对此深怀感激. 感谢康奈尔大学、北卡罗来纳州立大学和普渡大学的同事，尤其是Jim Berger、Larry Brown、David Cox爵士、Ziding Feng、Janet Johnson、Leon Gleser、Costas Goutis、Dave Lansky、George McCabe、Chuck McCulloch、Myra Samuels、Steve Schwager和Shayle Searle. 他们付出时间与专业知识，阅读本书部分手稿，提供帮助，并在多次讨论中提出建设性意见. 我们还感谢Shanti Gupta的热情接待，以及对本书至关重要的普渡大学图书馆. 感谢Shayle Searle，以及匿名和具名审稿人Jim Albert、Dan Coster、Tom Wehrly的细致阅读与有益建议. 感谢David Moore和George McCabe允许我们使用他们的表格，感谢Steve Hirdt提供数据. 由于两位作者在写书的大部分时间里相隔至少600英里，最后，我们感谢Bitnet，使这一切成为可能.
\begin{flushright}George Casella\\Roger L. Berger\end{flushright}
\clearpage
\begin{quote}
“我们已经谈到演绎和推断了，”雷斯垂德向我眨眨眼说，“福尔摩斯，我觉得光是处理事实就够难了，可不愿再飞到理论和幻想中去.”
\begin{flushright}雷斯垂德探长对夏洛克·福尔摩斯\\《博斯科姆比溪谷秘案》\end{flushright}
\end{quote}
